\begin{proofbox}

	\textbf{\textcolor{CProof}{思路提示}}

	证明原命题与逆否命题等价，可通过将蕴含式转化为析取式，再比较两者的逻辑公式是否相同。

	\medskip
	\textbf{\textcolor{CProof}{正式推导}}
	\begin{description}[style=nextline, leftmargin=2.8em, labelsep=0.5em, itemsep=0.55em]
		\item[\textbf{步骤一（蕴含转化）：}] 将原命题 $p\to q$ 用逻辑等价式表示：
			\[
				p\to q \equiv \lnot p \lor q.
			\]
			\par\textbf{理由：}
			蕴含式的真值定义：当 $p$ 为假或 $q$ 为真时 $p\to q$ 为真，正好对应 $\lnot p\lor q$。

		\item[\textbf{步骤二（逆否转化）：}] 将逆否命题 $\lnot q\to\lnot p$ 同样转化：
			\[
				\begin{aligned}
					\lnot q\to\lnot p &\equiv \lnot(\lnot q) \lor \lnot p \\
					&\equiv q \lor \lnot p.
			\end{aligned}
			\]
			\par\textbf{理由：}
			同样依据蕴含式转化，并将双重否定 $\lnot(\lnot q)$ 化简为 $q$。

		\item[\textbf{步骤三（对比等价）：}] 由析取交换律，$\lnot p\lor q$ 与 $q\lor\lnot p$ 是同一个公式，因此
			\[
				p\to q \equiv \lnot q\to\lnot p.
			\]
			\par\textbf{理由：}
			逻辑或满足交换律，因此 $\lnot p\lor q$ 和 $q\lor\lnot p$ 在逻辑上等价。
	\end{description}

	\medskip
	\textbf{\textcolor{CProof}{结论回扣}}

	\textbf{原命题与逆否命题恒等价，因此在推理中可以将原命题转化为逆否命题进行证明或判断。}
\end{proofbox}
